\section{信息系统与智能决策的参数影响及方法对照}\label{sec:rel-cyber-impact-cases}

信息系统不是一个可以直接乘在故障率上的总可用率。本章从共享元件、备选路径、服务质量、共因环境、物理供电、电池自治和
智能决策依次推导联合事件概率，并说明这些因素如何改变静态 FMEA、仅保护和信息物理联合三种结果。重点是保持同一故障
频率与同一后果基准，使差值能够归因。

\subsection{元件状态、路径状态与功能状态}
设信息元件集合为 $\mathcal C=\{1,\ldots,n\}$，元件 $i$ 的工作状态为 $x_i\in\{0,1\}$，固有可用率为 $a_i$。
在独立环境中，完整状态 $x$ 的概率为
\begin{equation}\label{eq:rel-cyber-state-probability}
 p(x)=\prod_{i=1}^{n}a_i^{x_i}(1-a_i)^{1-x_i}.
\end{equation}
一条路径 $r$ 是元件索引集合 $C_r$。路径结构可用当且仅当所有成员工作：
\begin{equation}
 X_r(x)=\prod_{i\in C_r}x_i.
\end{equation}
功能 $f$ 可以有若干备选路径 $\mathcal R_f$，其结构状态为
\begin{equation}
 X_f(x)=1-\prod_{r\in\mathcal R_f}(1-X_r(x)).
\end{equation}
平台不先把每条路径压缩成独立概率再做并联，因为两条路径可能共享交换机、电源或通信中心。它枚举元件位状态并对满足
$X_f=1$ 的状态概率求和，因此共享依赖只计算一次。

\subsection{串联、独立并联与共享并联的数值差异}
两元件串联且 $a_1=a_2=0.99$ 时，功能可用率为 $0.99^2=0.9801$。两条完全独立的单元件并联路径可用率为
\begin{equation}
 1-(1-0.99)^2=0.9999.
\end{equation}
若两条名义路径都经过可用率 $0.98$ 的共享中心，路径末端各为 $0.99$，则联合可用率不是把
$0.98\times0.99$ 当成两条独立路径并联，而是
\begin{equation}
 A_f=0.98\left[1-(1-0.99)^2\right]=0.979902.
\end{equation}
可见共享中心使“新增冗余链路”的收益从接近两个九降到受 $0.98$ 上限约束。程序的最小割集会识别共享中心单点割集和
两个末端同时失效的二元割集；这比只报告路径条数更能说明薄弱环节。

\subsection{QoS 是确定性准入，不是事后折扣}
每个信息元件给出时延 $d_i$、抖动 $j_i$ 和包交付概率 $q_i$。路径聚合为
\begin{align}
 d_r&=\sum_{i\in C_r}d_i,\\
 j_r&=\sum_{i\in C_r}j_i,\\
 q_r&=\prod_{i\in C_r}q_i.
\end{align}
功能契约给出 $d_f^{max}$、$j_f^{max}$、$q_f^{min}$。路径必须同时满足
\begin{equation}
 d_r\le d_f^{max},\qquad j_r\le j_f^{max},\qquad q_r\ge q_f^{min},
\end{equation}
才能参与结构可用性。一个固有可用率为 $0.999999$ 但保护消息时延 200 ms、契约上限 20 ms 的路径，对该保护功能的贡献
仍为零。把它乘成 $0.999999\times20/200$ 会制造没有物理依据的部分成功状态。

QoS 改善的影响具有门槛性。路径时延从 25 ms 降到 21 ms 时，20 ms 契约下功能可用率不变；降到 19 ms 时，整条路径
进入可用集合，功能可用率发生离散跃迁。因此参数扫描应同时报告准入路径集合，而不能只画一条平滑灵敏度曲线。

\subsection{共因环境的条件枚举}
设环境状态 $g\in\mathcal G$ 的概率为 $p_g$，环境指定失效共因组和失电物理母线。条件于 $g$ 后，属于失效共因组的元件
状态强制为零；其余元件再按式~\eqref{eq:rel-cyber-state-probability} 枚举。联合概率为
\begin{equation}
 p(g,x)=p_g p(x\mid g),\qquad
 \sum_g\sum_xp(g,x)=1.
\end{equation}
若正常环境概率为 $0.98$、通信中心共因失效环境为 $0.02$，即使所有独立元件可用率趋近于一，依赖该中心的功能可用率也
不会超过 $0.98$。这解释了独立元件升级后的收益饱和。

环境概率必须构成完备分布。若用户只给出若干异常环境且概率和小于一，入口不会暗自归一化并改变事件频率；调用方应显式
加入正常环境。概率为负、和超出容差或共因组引用不存在均拒绝执行。

\subsection{信息设备供电与电池自治}
信息元件可以绑定物理母线稳定 ID。若环境使该母线失电，元件仅在备用电池能够覆盖信息停电持续时间时继续工作。设电池
能量为 $E_b$ Wh、功耗为 $P_b$ W，则自治时间
\begin{equation}
 T_b=\begin{cases}E_b/P_b,&P_b>0,\\+\infty,&P_b=0.\end{cases}
\end{equation}
静态事件树中，若 $T_b\ge T_{out}$，该次环境下供电依赖不使元件失效；否则元件在整个代表事件中视为不可用。序贯队列
则在 $t_{deplete}=t_{out,start}+T_b$ 产生电池耗尽事件，因此同一次物理停电的前半段和后半段可对应不同的信息功能状态。

例如 $E_b=20$ Wh、$P_b=10$ W，自治时间为 2 h。面对 1.5 h 停电，静态与序贯模型都保持功能；面对 3 h 停电，静态
模型把功能判为失效，序贯模型保留前 2 h 自动操作能力，最后 1 h 进入人工或失败类。若关键控制动作发生在停电后 10 min，
序贯结果会优于静态二值近似。

\subsection{联合功能而非边际功能的必要性}
检测、主跳闸、后备跳闸、隔离和恢复可能依赖同一信息元件。分别计算五个边际可用率再相乘隐含了功能独立，通常会低估共享
中心失效的相关性。平台对同一个状态 $x$ 同时计算五个功能布尔值，得到功能向量
\begin{equation}
 y(x)=(y_D,y_{TP},y_{TB},y_I,y_R),
\end{equation}
再按相同向量归并概率。若一个共享中心控制所有功能，只有“全部可用”和“全部不可用”两个主要类；把五个相同边际概率
相乘会错误地产生大量部分成功概率。

对故障场景 $e$，保护信息联合类概率可写为
\begin{equation}
 \pi(c\mid e)=\sum_{g,x}p(g,x)\,
 \mathbf1\{\Phi(e,g,x)=c\},
\end{equation}
其中 $\Phi$ 顺序执行检测、主后备保护、隔离、恢复与 DER-FRT 分类。类概率和必须为一；概率为零的类不参与后果聚合。

\subsection{检测混淆矩阵}
真实故障类别为 $z$，检测输出为 $\widehat z$，混淆矩阵元素
$C_{z\widehat z}=P(\widehat z\mid z)$。每个真实类别行必须非负且和为一。检测结果还带有时延类 $\ell$，其概率为
$P(\ell\mid z,\widehat z)$。检测类的完整概率为
\begin{equation}
 p(z,\widehat z,\ell)=p(z)C_{z\widehat z}P(\ell\mid z,\widehat z).
\end{equation}
漏检进入检测失败类；错误分类可以进入错误隔离区，并不等同于漏检。提高总检出率若同时把概率从正确分类移向错误分类，可能
增加切负荷，因此“召回率更高”不能单独证明可靠性改善。

检测时延会扩展故障未隔离窗口。若该窗口切负荷为 $S_d$，时延从 $t_a$ 增至 $t_b$ 的 EENS 增量为
\begin{equation}
 \Delta E_d=\lambda p(z,\widehat z,\ell)S_d
 \frac{t_b-t_a}{3600}.
\end{equation}
秒到小时的单位换算必须在此处完成。对 10 MW、2 次/年、概率 0.5 的类别，时延增加 0.2 s 只增加
$5.56\times10^{-4}$ MWh/年；若程序显示 2 MWh/年，通常是遗漏了 $3600$。

\subsection{最小风险隔离}
候选隔离动作 $a$ 具有切除负荷 $S_a$、误隔离代价 $C_a^{false}$、漏隔离风险 $C_a^{miss}$ 和安全约束。给定检测后验
$b_z=P(z\mid\widehat z)$，动作风险为
\begin{equation}
 R(a\mid b)=C_a^{op}+\sum_zb_z C(a,z),
\end{equation}
其中 $C_a^{op}$ 可包含切负荷能量和开关动作成本。算法先删除引用不存在、无法覆盖故障或违反安全门槛的动作，再在可行集合
中选最小风险动作。若没有可行动作，结果显式为失败，不返回默认第一项。

考虑后验 $b=(0.8,0.2)$，窄区动作对类别 1 的代价为 1、对类别 2 的漏隔离代价为 20，宽区动作对两类代价均为 5，则
\begin{align}
 R_{narrow}&=0.8\times1+0.2\times20=4.8,\\
 R_{wide}&=0.8\times5+0.2\times5=5.0.
\end{align}
窄区略优；若类别 2 后验升到 0.25，窄区风险变为 5.75，算法改选宽区。这个切换点来自后验和代价，而不是写死的置信度
阈值。

\subsection{风险约束恢复}
恢复动作 $a$ 给出恢复负荷 $L_a$、不安全概率 $p_a^{unsafe}$、执行成功率 $p_a^{exec}$ 和期望剩余损失。平台先施加
\begin{equation}
 p_a^{unsafe}\le\bar p_{unsafe}
\end{equation}
硬约束，再在可行动作中最小化风险调整损失或最大化风险调整恢复量。一个恢复 10 MW 但不安全概率 0.2 的动作，在门槛
0.05 下不会因收益高而被选中。收紧门槛可能提高短期 EENS，却减少二次事故暴露；可靠性和安全性不能压成一个没有单位说明
的加权分数。

若动作恢复 $L_a$，执行失败时仍损失 $S_f$，成功后损失 $S_s$，则条件期望损失
\begin{equation}
 \overline S_a=p_a^{exec}S_s+(1-p_a^{exec})S_f.
\end{equation}
信息系统影响 $p_a^{exec}$ 时，应改变该条件期望或类概率；不能只在结果上附加“智能恢复已启用”布尔值。

\subsection{有限 POMDP 精确信念树}
部分可观测恢复模型用状态 $s\in\mathcal S$、动作 $a\in\mathcal A$、观测 $o\in\mathcal O$、转移
$T(s'\mid s,a)$、观测模型 $O(o\mid s',a)$ 和阶段代价 $c(s,a)$。给定信念 $b(s)$，观测后的贝叶斯更新为
\begin{equation}
 b'(s')=\frac{O(o\mid s',a)\sum_sT(s'\mid s,a)b(s)}
 {\sum_{\bar s}O(o\mid\bar s,a)\sum_sT(\bar s\mid s,a)b(s)}.
\end{equation}
有限时域价值递推为
\begin{equation}
 V_h(b)=\min_a\left[\sum_sb(s)c(s,a)+
 \sum_oP(o\mid b,a)V_{h-1}(\tau(b,a,o))\right].
\end{equation}
实现遍历完整有限信念树，不用固定“先测量后动作”规则。专项算例构造一个立即恢复风险高、先获取信息后可按状态选择动作的
系统，精确解必须选择信息收集动作。非法转移概率、观测概率或维数不一致直接拒绝。

\subsection{连续时间序贯事件队列}
信息元件 $i$ 的失效率、修复率为 $\lambda_i^c,\mu_i^c$，其两状态 CTMC 的下一事件时间分别服从指数分布。物理故障、信息
失效、修复、共因开始/结束、包投递、母线失电/复电和电池耗尽都进入按时间排序的同一队列。处理同一时刻事件时采用确定性
优先级，避免容器迭代顺序改变结果。

每次物理故障到达时，算法读取该时刻的信息状态和剩余电池能量，生成保护、隔离和恢复后果区间。年度 LOLE 不是简单相加
所有区间长度，而是对区间并集求测度：
\begin{equation}
 \mathrm{LOLE}_{y}=\left|\bigcup_{k\in\mathcal O_y}
 [t_k^{start},t_k^{end})\right|.
\end{equation}
EENS 对同时停电区间的切负荷可按系统约定聚合，但绝不能让同一 1 h 停电因两个重叠故障变成 2 h LOLE。专项测试构造
重叠事件并检查并集长度上界。

\subsection{包交付抽样与结构状态的区别}
元件结构工作不代表本次报文必达。对已通过 QoS 门槛的路径，包交付仍按 $q_r$ 抽样；若有多条可用路径，功能成功是至少一条
报文成功。结构失效是持续状态，包丢失是按消息发生的随机事件。把包交付概率并入元件长期可用率会错误地产生修复过程，
把结构失效只按单包抽样又会丢失持续相关性。

序贯结果返回包交付是否被抽样、信息供电耦合是否生效、电池能量轨迹是否生效和最小剩余能量。只有这些执行证据为真时，
才能声明相应因素进入年度指标。

\subsection{三级方法的共同基准}
三级对照的静态行、仅保护行和联合行必须共享：场景 ID、起始频率 $\lambda_e$、主/后备/未清除切负荷、隔离后切负荷、恢复后
切负荷、自动/人工恢复时间和修复时间。静态行把未清除切负荷保持到修复：
\begin{equation}
 E_e^S=S_{e,U}\tau_{e,repair}.
\end{equation}
仅保护行删除全部信息元件和功能绑定，只保留保护类：
\begin{equation}
 E_e^P=\sum_{c_P}\pi_P(c_P\mid e)W(e,c_P).
\end{equation}
联合行保留完整信息状态：
\begin{equation}
 E_e^C=\sum_{g,x,c_P}p(g,x)\pi(c_P\mid e,g,x)W(e,g,x,c_P).
\end{equation}
最终三行分别乘同一个 $\lambda_e$ 并求和。

\subsection{保护收益和信息代价}
平台定义
\begin{align}
 \Delta E_{prot}&=\mathrm{EENS}^{S}-\mathrm{EENS}^{P},\\
 \Delta E_{cyber}&=\mathrm{EENS}^{C}-\mathrm{EENS}^{P}.
\end{align}
$\Delta E_{prot}$ 表示相对静态未清除基线的保护收益；$\Delta E_{cyber}$ 表示相对理想信息条件的联合增量。二者不是 Shapley
分解，也不含无故障误动，不能相加后宣称得到全部信息物理贡献。

在后果单调且信息仅会使功能从成功变失败时，$\Delta E_{cyber}\ge0$。若远方自动控制错误会造成误切，而信息失效反而阻止
该错误动作，则差值可以为负。程序保留实际符号，不把负值截零；报告必须解释对应类的后果。

\subsection{闭式三级对照算例}
基准场景取年频率 2 次，未清除切负荷 10 MW、修复 10 h，因此静态行
\begin{equation}
 \mathrm{EENS}^{S}=2\times10\times10=200\ \text{MWh/年},
 \quad \mathrm{LOLE}^{S}=20\ \text{h/年}.
\end{equation}
主保护成功率 0.9、后备成功率 1，主/后备清除延时分别为 0.1 s 和 0.5 s；隔离、恢复后果使用同一三窗口。闭式事件树得到
\begin{align}
 \mathrm{EENS}^{P}&=0.4007777778\ \text{MWh/年},\\
 \mathrm{LOLE}^{P}&=0.2000777778\ \text{h/年}.
\end{align}
共享信息中心可用率为 0.8，且检测、跳闸、隔离、恢复都绑定该中心时，联合结果为
\begin{align}
 \mathrm{EENS}^{C}&=21.5217333333\ \text{MWh/年},\\
 \mathrm{LOLE}^{C}&=4.1600622222\ \text{h/年}.
\end{align}
因此保护收益为 $199.5992222222$ MWh/年，信息失效相对仅保护的增量为 $21.1209555555$ MWh/年。这个结果同时证明保护
逻辑和信息状态都实际改变了后果；若三行完全相同，说明运行链没有消费相应模型。

\subsection{信息可用率扫描}
保持静态与保护输入不变，把共享中心可用率从 0.8 降到 0.4。静态行必须严格保持 200 MWh/年，保护行也必须保持
$0.4007777778$ MWh/年；只有联合行上升。该测试比“联合结果非零”更有辨识力，因为它同时检查比较基准未被信息参数污染。

若有两条独立路径，单一路径可用率从 0.8 降到 0.4 对功能可用率的影响取决于另一条路径；若两条路径共享中心，中心参数
的边际影响最大。参数扫描必须与最小割集一起解释，不能只比较一个总百分比。

\subsection{有限差分敏感性检查}
对任一连续概率参数 $\theta$，可用中心差分
\begin{equation}
 D_h(\theta)=\frac{E(\theta+h)-E(\theta-h)}{2h}
\end{equation}
检查解析方向。$h$ 应避免跨越 $[0,1]$ 边界和 QoS 离散门槛。对纯串联独立路径，解析导数可直接由乘积得到；对共享依赖、
共因和功能绑定，精确状态枚举提供无抽样噪声的差分基准。序贯模型有随机误差时，应使用公共随机数或置信区间比较，不能以
单次样本差的符号判定错误。

\subsection{信息物理验收条件}
信息系统相关实现至少满足：元件状态概率和为一；共享路径与直接位枚举一致；QoS 超限路径不进入可用集合；最小割集确实使
功能失败且删除任一成员后不再是割集；共因环境条件概率守恒；母线失电和电池耗尽改变后续事件结果；包投递抽样标志为真；
检测混淆矩阵逐行守恒；隔离与恢复只在可行动作中选择；POMDP 与小规模手算价值一致；停电区间使用并集；三级对照的静态行
不随信息参数变化而联合行按类后果变化。实现分别位于
\srcpath{src/reliability/reliability\_assessment.cpp}、
\srcpath{src/reliability/intelligent\_cyber\_physical.cpp} 和
\srcpath{src/reliability/protection\_frt.cpp}，HTTP 与 GUI 的三级结果使用同一 C++ 聚合入口。
